% !PRINTED 270
\weberpage{270}
%<<CONT
so ergiebt sich, wenn $U = \begin{pmatrix}\lambda, & \mu\\ \nu, & \varrho\end{pmatrix}$ gesetzt wird, als Bedingung
\[
\begin{aligned}
\sigma^r\lambda &= \pm\sigma^t\lambda, & \sigma^r\mu &= \pm\sigma^{-t}\mu\\
\sigma^{-r}\nu &= \pm\sigma^t\nu, & \sigma^{-r}\varrho &= \pm\sigma^{-t}\varrho,
\end{aligned}
\]
wo überall dasselbe Zeichen gelten muss. Daraus folgt, dass entweder
\[
\sigma^r = \pm\sigma^t, \quad \nu = 0, \quad \mu = 0, \quad \lambda\varrho = 1
\]
oder
\[
\sigma^r = \pm\sigma^{-t}, \quad \lambda = 0, \quad \varrho = 0, \quad \mu\nu = -1
\]
sein muss, und dann ergiebt sich:
\[
USU^{-1} = S \quad \text{oder} \quad = S^{-1},
\]
und in beiden Fällen stimmen also die beiden Cyklen (12) vollständig mit einander überein.

Hiernach ergiebt sich aus den Zahlen am Schluss des \S.\,68 für die Anzahl der Cyklen erster, zweiter und dritter Art:
\[
p+1, \quad \frac{p(p+1)}{2}, \quad \frac{p(p-1)}{2}.
\]

Es ist zweckmässig, neben den Gruppen $L_p$ und $\Lambda_p$ noch eine dritte Gruppe in $E_{2,p}$ zu betrachten, die mit diesen isomorph ist, wenn sie sich auch nicht durch Transformation darauf zurückführen lässt.

Da $N$ reell ist, so können wir nach \S.\,67 eine Zahl $\nu$ in $\mathfrak{E}_{2,p}$ so bestimmen, dass $N = \nu^{p+1}$ ist. Wir lassen nun der Substitution
\[
A = \begin{pmatrix}\alpha, & \beta\\ -N\beta^p, & \alpha^p\end{pmatrix}
\]
aus $\Lambda_p$ eine Substitution $B$ entsprechen, die so gebildet ist:
\weq{13}{B = \begin{pmatrix}\alpha, & \nu\beta\\ -\nu^p\beta^p, & \alpha^p\end{pmatrix}, \quad \alpha^{p+1} + N\beta^{p+1} = 1.}
Setzen wir zwei Substitutionen $B$, $B_1$ von der Form (13) zusammen, so erhalten wir
\[
\begin{aligned}
BB_1 &= \begin{pmatrix}\alpha, & \nu\beta\\ -\nu^p\beta^p, & \alpha^p\end{pmatrix}\begin{pmatrix}\alpha_1, & \nu\beta_1\\ -\nu^p\beta_1^p, & \alpha_1^p\end{pmatrix}\\
&= \begin{pmatrix}\alpha\alpha_1 - N\beta\beta_1^p, & \nu(\alpha\beta_1+\beta\alpha_1^p)\\ -\nu^p(\beta^p\alpha_1+\alpha^p\beta_1^p), & \alpha^p\alpha_1^p - N\beta^p\beta_1\end{pmatrix};
\end{aligned}
\]
darin sind

$\alpha\alpha_1 - N\beta\beta_1^p, \ \alpha^p\alpha_1^p - N\beta^p\beta_1$ und $\alpha\beta_1+\beta\alpha_1^p, \ \beta^p\alpha_1+\alpha^p\beta_1^p$
%CONT>>
